\begin{abstract}
Este trabalho avalia a aplicação de modelos de aprendizado de máquina para a previsão
de demanda mensal em uma empresa de sucos naturais operando sob sistema de produção
puxada. Foram comparados um modelo de regressão linear múltipla e uma rede neural
recorrente do tipo LSTM, treinados sobre uma base de dados mensal de 15 sabores
comercializados entre abril de 2025 e junho de 2026. Os resultados, avaliados por
coeficiente de determinação (R²) tanto de forma agregada quanto individualizada por
sabor, indicam que a regressão linear apresenta desempenho mais estável do que a
LSTM na maioria dos sabores analisados, sobretudo quando a avaliação agregada — que
tende a mascarar falhas de ajuste em sabores de menor volume — é substituída por uma
avaliação desagregada. A partir dos resultados, propõe-se uma estratégia de gestão de
estoque baseada em três cenários de risco (conservador, de equilíbrio e otimista),
fundamentada no erro histórico de previsão obtido por validação cruzada. O estudo
discute ainda os desafios impostos por uma base de dados de tamanho reduzido e por
sabores de lançamento recente, oferecendo subsídios práticos para o planejamento
produtivo da empresa estudada.

\textbf{Palavras-chave:} previsão de demanda; aprendizado de máquina; regressão linear;
LSTM; produção puxada; gestão de estoque.
\end{abstract}